
\begin{figure}[htbp]
  \centering
  \includegraphics[ width=\linewidth]{figures/LZ_SR1_limit_SI_withFull.pdf}
  \caption{The \SI{90}{\percent} confidence limit (black line) for the spin-independent WIMP cross section vs.~WIMP mass. The gray dot-dash line shows the limit before applying the power constraint described in the text. The green and yellow bands are the $1\sigma$ and $2\sigma$ sensitivity bands. The dotted line shows the median of the sensitivity projection. Also shown are the PandaX-4T~\cite{PhysRevLett.127.261802}, XENON1T~\cite{collaboration2018dark}, LUX~\cite{akerib2017results}, and DEAP-3600~\cite{DEAP:2019yzn} limits. }
  \label{fig:limitplot}
\end{figure}

\begin{figure}[htbp]
  \centering
  \includegraphics[ width=\linewidth]{figures/model_fit_reconstructed_energy_lin_bins_revised.pdf}
  \caption{Reconstructed energy spectrum of the best-fit model. Data points are shown in black. The blue line shows total summed background. The darker blue band shows the model uncertainty and the lighter blue band the combined model and statistical uncertainty. Background components are shown in colors as given in the legend. Background components from $^8$B solar neutrinos and accidentals are included in the fit but are too small to be visible in the plot.}
  \label{fig:energy}
\end{figure}



Statistical inference of WIMP scattering cross section and mass is performed with an extended unbinned profile likelihood statistic in the $\log_{10}$\stwoc-\sonec\ observable space, with a two-sided construction of the \SI{90}{\percent} confidence bounds~\cite{baxter2021recommended}.
Background and signal component shapes are modeled in the observable space using the \textsc{geant4}-based package \textsc{baccarat}~\cite{akerib2021simulations,ALLISON2016186} and a custom simulation of the LZ detector response using the tuned \textsc{nest} model. The background component uncertainties are included as constraint terms in a combined fit of the background model to the data, the result of which is also shown in Table~\ref{tab:bkgds}. 

Above the smallest tested WIMP mass of \SI{9}{GeV/c\squared}, the best-fit number of WIMP events is zero, and the data are thus consistent with the background-only hypothesis.  Figure~\ref{fig:limitplot} shows the \SI{90}{\percent} confidence level upper limit on the spin-independent WIMP-nucleon cross section $\sigma_{\mathrm{SI}}$ as a function of mass. For WIMP masses between \SI{13}{GeV/c\squared} and \SI{36}{GeV/c\squared}, background fluctuations produce a limit that is substantially smaller than the median expected limit, as shown by the dot-dashed line in Fig.~\ref{fig:limitplot}. For these masses, the limit is constrained to a cross section such that the power of the alternate hypothesis is $\pi_{\rm crit}=0.16$~\cite{cowan2011power}.  This restricts the fluctuation to $1\sigma$ below the median expected limit. The introduction of the power constraint also introduces overcoverage, i.e., the coverage of the limit with the power constraint is greater than 90\%.
%The minimum of the limit curve is at $m_{\chi}=$~\SI{36}{GeV/c\squared} with a limit of $\sigma_{\mathrm{SI}}=$ \SI{9.2E-48}{cm\squared}.
The minimum of the limit curve is $\sigma_{\mathrm{SI}}=$ \SI{9.2E-48}{cm\squared} at $m_{\chi}=$~\SI{36}{GeV/c\squared}. The minimum of the unconstrained limit curve is \SI{6.2E-48}{cm\squared} at \SI{26}{GeV/c\squared}, and the minimum of the median expected limit is \SI{1.9E-47}{cm\squared} at \SI{43}{GeV/c\squared}.
The background model and data as a function of reconstructed energy are shown in Fig.~\ref{fig:energy}, and the data agree with the background-only model with a p-value of 0.96. LZ also reports the most sensitive limit on spin-dependent neutron scattering, detailed in the Appendix. A data release for this result is in the Supplemental Material~\footnote{See Supplemental Material at \url{\supplementalurl} for a data release for information shown in Figs.~\ref{fig:NRacc},~\ref{fig:s1s2wimp},~\ref{fig:limitplot},~\ref{fig:limitplotSDn}, and~\ref{fig:limitplotSDp}}.


